\section{Limitations and Ethical Considerations}
\label{sec:limits-ethics}

\textbf{Limitations.} We acknowledge several limitations. (1)~\emph{Agent harness.} We use a minimal ReAct harness as the primary configuration; however, our harness ablation (\Cref{sec:res-harness}) shows that silent failures persist under single-shot and few-shot variants, indicating they are not an artifact of the minimal ReAct design. Production agents (e.g., code-interpreter, AutoGen) may still exhibit different failure distributions, and testing richer harnesses is future work. (2)~\emph{Model scope.} DeepSeek-R1 and Qwen3-max achieved 0\% success on UCI tasks, likely due to output-format incompatibility with the ReAct parser rather than fundamental inability. These models should not be called ``weakest'' but rather ``lowest-performing under the current harness.'' (3)~\emph{Taxonomy validation.} Cohen's $\kappa = 0.435$ (moderate) between rule-based and LLM-judged labels; the remaining disagreement centers on the boundary between \texttt{output\_mismatch} and \texttt{answer\_error}. Full human annotation with $\geq 3$ annotators is needed for a definitive $\kappa$. (4)~\emph{Bifurcation.} The null-model analysis shows $R^2 = 0.493$, reported as a descriptive statistic; the bifurcation is partially mechanical (SFR and execution-success share the loud-failure count) and cannot be decomposed into ``mechanical'' vs.\ ``real'' components. (5)~\emph{Oracle repair.} The 64\% repair rate matches the no-op baseline, demonstrating reparability but not causal attribution; full necessity/sufficiency analysis with human root-cause labels is future work. (6)~\emph{Statistical tests.} We use McNemar tests for binary verifier comparisons; multinomial mixed-effects models with task/model random effects would be more powerful and are left to future work.

\textbf{Ethical considerations.} This work analyses LLM agent failures on data-science tasks. No human subjects or personal data are involved. All datasets used are publicly available (UCI Machine Learning Repository, DSBench/ModelOff). The API costs (\$80 total) were modest. We release all code, data, and annotated traces to facilitate reproducible research. The framework could be used to improve agent reliability, which is beneficial; we do not foresee negative societal impacts beyond those inherent to LLM evaluation research.

\textbf{Generative AI usage.} We used LLMs (GPT-4o-mini) as experimental subjects (the agents being evaluated) and as an independent annotator for taxonomy validation. All LLM-generated outputs are treated as data, not as authorial content. The paper itself was written by the authors.
